\section{Context Compression Techniques}
\label{sec:compression}

Context compression techniques address the fundamental challenge of efficiently managing long contexts in large language models while maintaining performance. As modern applications increasingly demand the processing of extensive documents, lengthy conversations, and complex reasoning chains, the need to compress contextual information without sacrificing semantic integrity has become paramount. The field has evolved from simple token removal strategies to sophisticated learned representations that capture semantic content in compressed form, reflecting diverse trade-offs between compression ratio, computational cost, and performance preservation.

The landscape of context compression can be understood through multiple complementary perspectives. At the highest level, methods divide into soft prompt compression, which employs learned continuous representations to encode contextual information into compact embeddings, and hard prompt compression, which explicitly removes or reorders tokens to reduce sequence length. Beyond this dichotomy, architectural approaches modify the underlying transformer structure to enable efficient memory management, while information-theoretic frameworks establish fundamental bounds on what compression can achieve. Recent comprehensive surveys \citep{li2025prompt} have identified additional perspectives including attention optimization, parameter-efficient fine-tuning, modality integration, and synthetic language approaches, each offering unique insights into compression mechanisms.

\subsection{Soft Prompt Compression Through Learned Representations}

Soft prompt compression methods transform contextual information into dense vector representations that serve as compact substitutes for lengthy text sequences. These approaches leverage the representational capacity of neural networks to encode semantic information in continuous embedding spaces, enabling significant compression while maintaining task-relevant information through learned transformations.

\subsubsection{Gisting and Representational Bottlenecks}

The gisting approach introduced by \citet{mu2023learning} represents a conceptually elegant solution to prompt compression through the introduction of representational bottlenecks in attention mechanisms. By modifying standard transformer attention masks to force models to condense input information into dedicated ``gist'' token activations, this method creates context-specific compression in a single forward pass. The approach requires minimal architectural modifications, introducing only new vocabulary token embeddings and attention mask adjustments, making it readily integrable with existing frameworks and systems.

Gisting achieves up to 26$\times$ compression on LLaMA-7B and FLAN-T5-XXL models, yielding up to 40\% FLOPs reductions and 4.2\% wall time speedups with minimal quality loss on tasks amenable to high compression ratios. The method's training efficiency stands as a particular advantage, as gist models can be trained with no additional cost beyond standard instruction fine-tuning through straightforward attention mask modifications. However, performance degrades rapidly when compressing longer contexts beyond hundreds of tokens, limiting applicability to contemporary long-context scenarios. The fundamental limitation stems from the fixed-capacity bottleneck, which cannot effectively scale to the diverse information densities present in extended sequences. Despite these constraints, gisting demonstrates that architectural modifications to attention mechanisms can enable compression without explicit token removal or summarization.

\subsubsection{AutoCompressors and Summary Vector Generation}

\citet{chevalier2023adapting} advanced soft prompt compression through AutoCompressors, which adapt pre-trained language models into specialized compression models capable of transforming long contexts into summary vectors accessible as soft prompts. Unlike gisting's attention-based bottleneck, AutoCompressors employ an unsupervised training objective where long documents are processed in segments, with summary vectors from all previous segments incorporated into language modeling of subsequent segments. This segmented approach enables processing of sequences up to 30,720 tokens during fine-tuning of OPT and Llama-2 models.

The AutoCompressor framework demonstrates that summary vectors serve as effective substitutes for plain-text demonstrations in in-context learning scenarios, simultaneously improving perplexity on long-context tasks while reducing inference costs. The method provides a simple and inexpensive solution to context window extension, requiring only standard fine-tuning procedures without architectural modifications or complex training objectives. By training on massive text data with language modeling objectives, AutoCompressors learn to extract and preserve task-relevant information in compact vector representations. The learned summary vectors maintain semantic information effectively, as evidenced by their ability to substitute for full-text demonstrations in few-shot learning scenarios. This capability suggests that the compression process captures abstract task-relevant patterns rather than merely memorizing surface-level token sequences.

The success of AutoCompressors highlights the potential of unsupervised learning for compression, where models discover compression strategies through exposure to diverse text without explicit supervision on what information to retain or discard. This approach contrasts with supervised methods that require labeled data indicating important content, instead relying on the implicit signal provided by language modeling objectives to guide compression decisions.

\subsubsection{In-Context Autoencoder Architecture}

\citet{ge2024icae} proposed the In-Context Autoencoder (ICAE), which frames context compression as an autoencoding problem with a learnable encoder and fixed decoder. The encoder, adapted from a language model using Low-Rank Adaptation (LoRA), compresses long contexts into a limited number of memory slots, while the decoder (the target LLM) conditions on these compressed representations for various downstream purposes. This architectural separation enables efficient training and deployment, as the encoder requires only approximately 1\% additional parameters while the decoder remains frozen.

ICAE employs a two-stage training process that first pre-trains the encoder with both autoencoding and language modeling objectives on massive text data, ensuring that memory slots accurately represent original context. The second stage fine-tunes the system on instruction data to produce desirable responses to diverse prompts. This staged approach separates the learning of compression from the learning of task-specific behaviors, potentially improving both objectives through focused optimization. The method achieves 4$\times$ context compression with improved inference latency and reduced GPU memory cost, demonstrating practical benefits for deployment scenarios.

The autoencoder framing provides theoretical grounding for compression through the lens of information bottlenecks and rate-distortion theory. By forcing information through a constrained number of memory slots, ICAE creates an explicit bottleneck that encourages retention of only essential information. The use of LoRA for encoder adaptation ensures parameter efficiency while maintaining model capacity for learning complex compression functions. The combination of autoencoding objectives (which encourage reconstruction fidelity) and language modeling objectives (which encourage semantic coherence) creates a multi-objective optimization that balances compression quality with downstream task performance.

\subsection{Hard Prompt Compression Through Token Manipulation}

Hard prompt compression methods directly modify prompt text through token removal, reordering, or paraphrasing, creating shorter sequences that preserve essential semantic content. These approaches offer high interpretability, as compressed prompts remain human-readable text, and require minimal or no additional training, enabling straightforward deployment with existing language models.

\subsubsection{LLMLingua and Coarse-to-Fine Compression}

\citet{jiang2023llmlingua} introduced LLMLingua, a coarse-to-fine prompt compression method that leverages small language models to identify and remove low-importance tokens. The approach employs perplexity from well-trained small models (GPT-2-small or LLaMA-7B) as a proxy for token importance, with higher perplexity indicating greater information content and lower perplexity suggesting redundancy or predictability. This perplexity-based ranking enables principled token selection without task-specific training.

The LLMLingua framework incorporates three key components that work synergistically to achieve high-quality compression. First, a budget controller maintains semantic integrity under high compression ratios by allocating compression budgets across prompt segments based on their estimated importance. Second, token-level iterative compression models interdependence between compressed contents, recognizing that token importance depends on which other tokens remain in the compressed prompt. Third, instruction tuning achieves distribution alignment between the small compression model and target LLM, ensuring that importance judgments transfer effectively across models with different training distributions.

The coarse-to-fine strategy operates hierarchically, first identifying important prompt segments at a document or paragraph level, then performing fine-grained token selection within retained segments. This multi-scale approach proves more effective than uniform token pruning, as it recognizes that information importance varies both within and between segments. LLMLingua achieves up to 20$\times$ compression with minimal performance loss across diverse benchmarks including GSM8K for mathematical reasoning, BBH for challenging reasoning tasks, ShareGPT for conversational contexts, and Arxiv-March23 for summarization. The method's state-of-the-art results demonstrate that carefully designed hard prompt compression can rival or exceed soft prompt methods while maintaining interpretability and requiring no architectural modifications.

\subsubsection{LongLLMLingua for Extended Contexts}

Building upon the LLMLingua foundation, \citet{jiang2024longllmlingua} developed LongLLMLingua to address specific challenges in long-context scenarios, including higher computational cost, performance reduction due to the ``lost in the middle'' phenomenon, and position bias in transformer attention. The method extends the original coarse-to-fine compression with four key innovations tailored to long-context processing.

The question-aware coarse-to-fine compression mechanism conditions compression decisions on the specific query or task, recognizing that information importance is inherently task-dependent. This query-aware approach enables more aggressive compression of query-irrelevant content while preserving query-relevant information at higher fidelity. The document reordering mechanism strategically positions documents within the context window to mitigate position bias, placing high-importance documents at the beginning and end where transformer models demonstrate better attention and retrieval. Dynamic compression ratios adapt to varying information densities across documents, applying aggressive compression to low-information content and conservative compression to information-dense segments. Finally, the subsequence recovery strategy improves the model's perception of key information by maintaining important multi-token sequences intact rather than fragmenting them through token-level pruning.

LongLLMLingua achieves remarkable performance improvements, including up to 21.4\% accuracy gains on NaturalQuestions with 4$\times$ compression, 94\% cost reduction on the LooGLE long-context benchmark, and 1.4$\times$ to 2.6$\times$ latency acceleration on 10,000-token prompts with 2$\times$ to 6$\times$ compression ratios. These results demonstrate that compression can not only reduce costs but actually improve task performance by reducing noise and mitigating position bias effects. The method's effectiveness across chain-of-thought reasoning, long-context comprehension, and retrieval-augmented generation scenarios highlights its broad applicability to diverse long-context applications.

\subsubsection{Selective Context and Redundancy Reduction}

\citet{li2023compressing} proposed Selective Context, which identifies and prunes redundancy in input context to create more compact representations. Unlike LLMLingua's perplexity-based approach, Selective Context focuses explicitly on redundancy detection, seeking to remove repeated or highly predictable information while preserving unique content. The method achieves 50\% context cost reduction with 36\% memory savings and 32\% inference time reduction, while incurring minimal performance degradation of only 0.023 BERTscore and 0.038 faithfulness on evaluation metrics.

Selective Context demonstrates effectiveness across diverse document types including arXiv papers, news articles, and long conversations, and across tasks including summarization, question answering, and response generation. The method's task-agnostic nature enables deployment without task-specific tuning, providing a general-purpose compression solution applicable to varied scenarios. The conservative compression ratios (typically 2$\times$ to 4$\times$) prioritize performance preservation over aggressive size reduction, positioning Selective Context as a safe, production-ready compression approach for applications where quality maintenance is paramount.

The redundancy-focused perspective offers complementary insights to perplexity-based methods. While low perplexity indicates predictability from a language modeling perspective, redundancy detection identifies repetition and overlap across context segments. Content may exhibit low perplexity due to grammatical regularity or common phrasing while still conveying unique semantic information, whereas redundant content explicitly repeats previously stated information. By targeting redundancy specifically, Selective Context preserves diverse information even when that information follows predictable linguistic patterns.

\subsection{Advanced Hard Compression Techniques}

\paragraph{Dynamic Token Pruning}

Recent advances in token pruning have introduced dynamic selection mechanisms that adapt compression to generation context. LazyLLM, developed by Apple Machine Learning Research, enables language models to dynamically select different subsets of tokens from context in different generation steps, recognizing that token importance varies across the generation process. The method uses attention scores from prior transformer layers to measure token importance for predicting the next token, selectively computing key-value pairs for important tokens while deferring computation of remaining tokens to later steps when they may become relevant.

This lazy computation approach achieves efficient long-context processing by reducing computation while maintaining token availability, ensuring that all context tokens remain accessible when needed. Unlike static pruning methods that permanently remove tokens before generation, dynamic pruning adapts to the evolving information needs of generation, potentially accessing different context subsets for different generated tokens. This flexibility better aligns with the observation that different parts of long contexts become relevant at different points in generation, particularly for complex reasoning or multi-hop question answering.

\paragraph{Token Merging and Structured Pruning}

Beyond individual token pruning, recent work has explored token merging strategies that combine similar tokens in the key-value cache to reduce distinct entries while maintaining semantic coverage. Clustering-based approaches group tokens by embedding similarity, replacing clusters with representative tokens or merged representations. This approach proves particularly effective when context contains substantial redundancy or repeated concepts expressed through varied phrasing.

Structured pruning extends token-level compression to higher-level units including attention heads and entire layers. Research on attention head pruning has revealed that only a small subset of attention heads prove important for most tasks, with important heads exhibiting interpretable functions such as attending to adjacent words or tracking specific syntactic relations. The vast majority of heads, especially in encoder self-attention, can be removed with minimal performance impact. Methods employing A* search for head selection have demonstrated the ability to eliminate redundant attention heads while maintaining strict accuracy guarantees through optimal head subset selection.

\subsection{Architectural Memory Compression}

\subsubsection{Compressive Transformers}

\citet{rae2020compressive} introduced the Compressive Transformer, which extends the Transformer-XL memory mechanism by compressing old activations into compressed memory rather than discarding them when the attention window slides forward. This approach maintains short-term memory of recent activations in full detail while compacting older activations into lower-resolution compressed representations, creating a two-tier memory system that balances detail and coverage.

The compression mechanism employs convolutional pooling to aggregate older hidden states into compact representations, with hyperparameters controlling the compression ratio and the boundary between detailed recent memory and compressed historical memory. Attention mechanisms span both recent and compressed memories, enabling the model to access both fine-grained recent context and coarse-grained historical context. This design proves particularly effective for tasks requiring long-range dependencies where distant context provides useful signal even in compressed form.

Compressive Transformers achieved state-of-the-art results on WikiText-103 (17.1 perplexity) and Enwik8 (0.97 bits per character) language modeling benchmarks, demonstrating that compressed memories can effectively substitute for full-detail long-range context. The approach generalizes beyond language, proving effective for high-frequency speech modeling and serving as a memory mechanism for reinforcement learning agents on visual reasoning tasks. The introduction of the PG-19 benchmark, derived from books to provide challenging long-range modeling tasks, has promoted research in long-sequence learning.

The architectural integration of compression into the model structure offers advantages over post-hoc compression methods, as the model learns to leverage compressed representations during training rather than adapting to externally imposed compression. This joint learning enables the development of attention patterns and feature representations optimized for compressed memory access. However, architectural approaches require training models from scratch or extensive fine-tuning to incorporate compression mechanisms, whereas post-hoc methods can be applied to existing pre-trained models without retraining.

\subsubsection{Recurrent Context Compression}

\citet{ge2024recurrent} proposed Recurrent Context Compression (RCC), which extends LLM context windows through iterative segment compression inspired by state space models and recurrent neural networks. For sequences exceeding the maximum context length, RCC divides input into fixed-length segments and compresses each segment sequentially, generating state vectors that capture essential information. The concatenation of state vectors from all segments provides compressed historical context accessible during inference.

The recurrent structure proves particularly effective for extremely long sequences, demonstrating up to 32$\times$ compression on text reconstruction with BLEU4 scores near 0.95 and achieving nearly 100\% accuracy on passkey retrieval tasks with sequences up to one million tokens. These results demonstrate that iterative compression can maintain information fidelity across extreme sequence lengths when compression operates hierarchically across segments rather than attempting to compress entire sequences in a single step.

RCC addresses a unique challenge in instruction-context compression interference, where compressing both task instructions and context into the same format reduces instruction clarity and degrades model performance. The proposed instruction reconstruction strategy detects and recovers instruction quality while maintaining context compression, improving downstream task performance by separating instruction handling from context compression. This separation recognizes that instructions and context serve different roles, with instructions providing task specification requiring high fidelity and context providing information that may tolerate greater compression.

The recurrent compression architecture enables scalable processing of arbitrary-length sequences with linear memory growth through compression, contrasting with the quadratic memory requirements of standard transformer attention. The fixed-size state vectors provide constant memory overhead per segment regardless of segment length, enabling practical deployment on memory-constrained devices for extremely long sequences.

\subsection{Key-Value Cache Compression}

Key-value caching stores key and value matrices during inference to avoid recomputation, but cache memory grows linearly with sequence length, consuming substantial memory for long contexts. Each generated token requires storing key-value pairs for all previous tokens across all transformer layers and attention heads, leading to memory consumption that can exceed model parameter memory for long sequences.

\subsubsection{Token Eviction Strategies}

StreamingLLM addresses cache growth by maintaining key-value states for only the first few tokens (typically 4) combined with a sliding window of recent tokens. The initial tokens serve as ``attention sinks'' that anchor attention patterns, while recent tokens provide current generation context. This fixed-size cache enables processing of arbitrarily long sequences with constant memory, though at the cost of losing middle context information.

H2O (Heavy Hitter Oracle) dynamically identifies important tokens based on cumulative attention scores during generation, maintaining a two-part cache structure with budgets for recent tokens and ``heavy hitter'' tokens that receive high attention. This adaptive approach preserves tokens demonstrating high importance through empirical attention patterns while evicting tokens that receive minimal attention. The dynamic identification enables more effective compression than static rules, as importance emerges from actual model behavior rather than heuristic assumptions.

SnapKV compresses the key-value cache during the prefill stage using a small observation window at the prompt end to predict token importance. This early compression reduces cache size from the start of generation rather than waiting for cache growth to trigger eviction. The observation window principle recognizes that token importance for generation often becomes apparent from attention patterns in a limited context window, enabling predictive importance estimation without full-sequence processing.

\subsubsection{Principled KV Compression}

Recent work on Expected Attention \citep{zhang2025expected} provides a principled framework for key-value compression by estimating KV pair importance through predicting how future queries will attend to them. The method analyzes the expected distribution of future queries and computes expected attention scores in closed form, enabling ranking of key-value pairs by predicted future importance. This forward-looking approach proves more effective than backward-looking methods based on historical attention, as it optimizes for future utility rather than past usage.

The training-free nature of Expected Attention enables deployment without model modification or fine-tuning, relying instead on mathematical analysis of attention mechanisms and activation properties. The closed-form computation provides efficiency advantages over simulation-based importance estimation, which would require sampling multiple future generation trajectories to estimate expected attention.

Model-directed compression approaches enable models to indicate which information to discard through learned importance scoring, potentially adapting compression to task requirements and input characteristics. Adaptive KV cache compression methods like ACC-RAG dynamically adjust compression rates based on input complexity, avoiding over-compression of simple queries and under-compression of complex queries. This complexity-aware adaptation recognizes that optimal compression ratios vary substantially across queries, with simple factual questions tolerating aggressive compression while complex multi-hop reasoning requires more conservative compression to preserve reasoning chains.

\subsubsection{Compression Techniques and Limits}

Key-value compression employs diverse techniques beyond token selection. Key quantization reduces the precision of key vectors while maintaining semantic information, achieving significant size reduction with minimal accuracy loss through careful bit-width selection and quantization-aware training. KIVI demonstrated that 2-bit asymmetric quantization of key-value caches enables longer context generation on consumer GPUs without quality degradation, making long-context inference accessible on resource-constrained hardware.

Value offloading moves some value vectors to slower storage (CPU memory or disk), fetching them when needed for attention computation. This approach trades memory for latency, enabling processing of longer contexts than GPU memory alone permits while incurring access latency for offloaded values. Selective offloading strategies prioritize keeping high-importance values in fast memory while offloading low-importance values, minimizing latency impact.

Recent theoretical work has begun exploring fundamental limits of key-value cache compression, investigating whether compression can continue indefinitely or whether phase transitions exist beyond which further compression causes catastrophic performance degradation. Early findings suggest that attention mechanisms have inherent selectivity enabling substantial compression for many tasks, but information-theoretic bounds apply to compression limits for tasks requiring fine-grained context access.

\subsection{Context Compression for Retrieval-Augmented Generation}

Retrieval-augmented generation systems face unique compression challenges stemming from the need to process multiple retrieved documents with varying relevance, handle position bias effects in multi-document contexts, maintain factual grounding and citation capability, and balance information diversity with redundancy elimination.

\subsubsection{Traditional Contextual Compression}

Contextual compression in RAG iterates over retrieved documents, processing each through a language model to extract query-relevant information and returning only pertinent content. This approach removes irrelevant content and reduces noise, focusing on query-specific information and improving reasoning quality through noise reduction. The query-aware filtering proves more effective than query-agnostic compression, as relevance judgments depend fundamentally on information needs expressed in queries.

\subsubsection{Extreme Compression for RAG}

\citet{xrag2024} introduced xRAG (Extreme Context Compression for RAG), which achieves one-token-per-document compression by reinterpreting document embeddings as modality features and fusing them with LLM representation space. This modality fusion framework treats dense retrieval embeddings as a retrieval modality analogous to vision or audio modalities in multi-modal models, integrating them directly into LLM representations without textual intermediaries.

The approach maps retrieval embeddings into the LLM's representation space through learned projection layers, preserving semantic information while enabling the LLM to process compressed document representations. Each document compresses to a single token representation, enabling processing of dozens or hundreds of documents within limited context windows. xRAG achieves an average 10\% improvement across six knowledge-intensive tasks, significantly surpassing previous compression methods while dramatically reducing context requirements.

The extreme compression ratio raises questions about information preservation, as a single token must encode the content of entire documents that may contain thousands of tokens. The success of this approach suggests that for retrieval-augmented scenarios where documents serve primarily to provide factual information to answer specific queries, the relevant information may be remarkably compressible. The retrieval embedding already captures semantic content optimized for relevance matching, and the projection into LLM space enables the model to leverage this semantic encoding directly.

\subsubsection{Adaptive Compression for RAG}

Adaptive Context Compression for RAG (ACC-RAG) addresses the limitation of fixed compression rates by dynamically adjusting compression based on query complexity. Simple queries may require only basic facts from retrieved documents, tolerating aggressive compression, while complex multi-hop reasoning queries require preserving reasoning chains and connections across documents, necessitating conservative compression.

ACC-RAG analyzes query complexity through features such as query length, entity count, syntactic complexity, and semantic ambiguity, then assigns compression ratios adaptively. Document-specific compression enables different rates for different retrieved documents based on relevance scores and information density, recognizing that top-ranked documents merit more careful preservation than marginally relevant documents. The adaptive approach achieves better accuracy on complex queries through conservative compression while maintaining faster inference on simple queries through aggressive compression, optimizing the performance-efficiency trade-off across varied query distributions.

\subsubsection{RAG-Specific Compression Considerations}

Multi-document handling in RAG requires careful consideration of document interactions and dependencies, as answers may require synthesizing information across multiple documents. Naive compression of documents independently may break cross-document references or eliminate complementary information distributed across documents. Document reordering strategies mitigate the ``lost in the middle'' phenomenon by positioning high-relevance documents at context extremes where transformer models demonstrate better attention and information access.

Factual grounding and citation preservation present challenges for aggressive compression, as models must trace information to source documents for fact-checking and attribution. Token-level compression may fragment citations or separate factual claims from their sources, reducing groundedness and increasing hallucination risk. Hierarchical compression approaches that maintain document-level structure while compressing within documents can preserve source attribution while achieving significant compression.

\subsection{Semantic Preservation in Compression}

Semantic preservation represents a central challenge in context compression, as the goal extends beyond mere size reduction to maintaining meaningful information content. Lossy compression accepts some information loss in exchange for high compression ratios, proving appropriate for lengthy contexts with substantial redundancy, while lossless compression preserves all information but achieves minimal compression ratios and remains rarely achievable with LLMs due to linguistic redundancy and context-dependent semantics.

\subsubsection{Measurement and Evaluation}

Embedding similarity provides a direct measure of semantic preservation through comparing embeddings of original and compressed text, with high cosine similarity indicating successful preservation of semantic content. BERTscore and similar metrics quantify semantic similarity at token and phrase levels, capturing fine-grained preservation quality. Downstream task performance offers ultimate validation, measuring whether compressed context supports task completion as effectively as full context across question answering, summarization, classification, and information extraction tasks.

Information retention metrics including ROUGE scores for summarization, F1 scores for factual content, and faithfulness measures quantify specific aspects of preservation quality. These metrics capture whether compression maintains key facts, preserves relationships between concepts, and supports required reasoning patterns. Multi-dimensional evaluation across diverse metrics provides comprehensive assessment of compression quality, as different applications prioritize different preservation aspects.

\subsubsection{Advanced Preservation Methods}

ChunkKV employs semantic chunks as compression units rather than individual tokens, grouping tokens into semantic units and retaining the most informative chunks while discarding less important tokens. This chunk-based granularity maintains semantic relationships better than token-level pruning, as meaningful units remain intact rather than fragmenting across pruning decisions. The approach demonstrates that compression unit selection significantly impacts preservation quality, with semantically-motivated units outperforming arbitrary tokenization boundaries.

Hybrid Context Compression (HyCo2) integrates global and local perspectives, balancing essential semantics captured in document-level summaries with critical details requiring fine-grained preservation. The hybrid strategy adapts compression to content type, applying different methods to different content segments based on information characteristics. This flexibility enables tailoring compression to content properties, applying aggressive summarization to background information while preserving technical details or factual claims at higher fidelity.

AdmTree employs adaptive semantic trees as a hierarchical compression framework, dynamically segmenting content based on information density with variable-length compression units. The semantic tree construction creates hierarchical representations capturing relationships between concepts at multiple granularities. This structured compression preserves both individual facts and their relationships, supporting reasoning that requires understanding connections between concepts.

\subsubsection{Context-Aware Preservation Strategies}

Hierarchical preservation categorizes information into primary (core facts and concepts compressed minimally), secondary (supporting details compressed moderately), and tertiary (background information compressed aggressively) tiers. This tiered approach enables balanced compression that maintains essential information while aggressively compressing less critical content. The flexible trade-offs allow tuning preservation-compression ratios to application requirements, with fact-intensive applications prioritizing primary information preservation and summarization applications tolerating more aggressive compression.

Query-aware compression adjusts preservation based on query information needs, identifying query-relevant content for careful preservation while aggressively compressing query-irrelevant content. This task-specific optimization recognizes that information importance depends fundamentally on intended use, with the same document requiring different compression strategies for different queries. Gradient-based importance estimation quantifies content contribution to task loss, providing direct measurement of preservation necessity for downstream task performance.

\subsection{Information-Theoretic Perspectives}

Information theory provides fundamental limits on context compression through establishing connections between compression, entropy, and prediction. The Kolmogorov complexity framework reveals that optimal prediction of data sequences is intimately tied to most efficient compression, with the Kolmogorov structure function measuring optimal compressor performance under model-size constraints. This theoretical foundation illuminates compression limits and enables analysis of compression method optimality.

\paragraph{Compression Bounds and Effective Context}

Finite-capacity models must compress information, with compression necessarily introducing errors on incompressible data. Information-theoretic constraints bound attainable accuracy on decidable tasks, with finite description length enforcing compression error and long-tail factual knowledge requiring prohibitive sample complexity. These fundamental limits constrain what any compression method can achieve, establishing performance ceilings independent of algorithmic sophistication.

Even with extremely long nominal context windows such as 128,000 tokens, geometric and computational effects compress contexts far below nominal size. Positional under-training causes effective compression because models receive limited training on extreme positions, leading to degraded performance at positions rarely encountered during training. Encoding attenuation effects reduce information density as sinusoidal or rotary positional encodings decay with distance, causing attention scores to decay for distant positions. Softmax crowding creates numerical bottlenecks through information crowding in attention computation, with score margin compression at long distances concentrating entropy in recent positions.

These effects jointly cause effective context to scale sub-linearly with nominal length, with gradient decay at rare positions, vanishing overlap in positional encodings, and logarithmic score-margin growth limiting utilization. The practical implication is that even without explicit compression, long contexts experience substantial implicit compression through these architectural and optimization effects, suggesting that explicit compression methods that align with these implicit compression patterns may achieve better quality-ratio trade-offs.

\paragraph{Entropy Law and Compression-Performance Connections}

Recent work has identified an ``entropy law'' connecting LLM performance with data compression ratio and first-epoch training loss, reflecting the information redundancy of datasets. Language models output probability distributions over next tokens conditioned on context, information that enables near-optimal compression via arithmetic coding. Shannon's theorem guarantees optimality if probabilities match true distributions, establishing an elegant equivalence between language modeling and data compression where better language models yield better compressors.

This equivalence provides theoretical grounding for compression approaches, as language modeling objectives implicitly train models to identify and eliminate redundancy. The entropy of context bounds achievable compression ratio, with repetitive or predictable content enabling higher compression than information-dense technical content. Compression methods that leverage language model probabilities to guide compression decisions thus align with information-theoretic optimality, explaining the effectiveness of perplexity-based approaches like LLMLingua.

The rate-distortion framework formalizes trade-offs between compression ratio and distortion (information loss), establishing theoretical limits on useful compression ratios. With a frozen language model, hard limits exist for compression via token dropping or reordering, as any compression must discard information that the fixed model cannot reconstruct. Learned compression methods that adapt model representations to compressed inputs potentially achieve better rate-distortion trade-offs by enabling the model to leverage compression-specific features.

\subsection{Summary and Comparative Analysis}

Context compression techniques span a spectrum from aggressive methods achieving 20$\times$ to 32$\times$ compression through LLMLingua, Gisting, and Recurrent Context Compression, to conservative approaches achieving 2$\times$ to 4$\times$ compression through Selective Context and semantic compression methods, with intermediate methods such as ICAE and AutoCompressors balancing 4$\times$ to 8$\times$ compression with performance preservation. Training costs vary similarly, from zero additional cost for training-free methods like LLMLingua and Selective Context, through low costs for Gisting's attention mask modifications and ICAE's 1\% parameter overhead with LoRA, to moderate costs for AutoCompressors' fine-tuning on long sequences, and substantial costs for Compressive Transformers requiring architecture-specific training from scratch.

Application domains reveal method specialization, with RAG systems benefiting particularly from LongLLMLingua and xRAG's document-handling capabilities, in-context learning scenarios favoring Gisting and AutoCompressors for demonstration compression, long document processing leveraging ICAE and Compressive Transformers' architectural approaches, and general-purpose applications employing semantic compression and Selective Context for broad applicability. The diversity of approaches reflects fundamental trade-offs between compression ratio (higher ratios reduce costs but risk information loss), computational cost (training-free methods enable immediate deployment while learned methods may achieve better quality), performance preservation (conservative compression maintains quality while aggressive compression maximizes efficiency), and interpretability (hard prompt methods maintain human readability while soft prompt methods achieve better compression through learned representations).

Recent comprehensive surveys \citep{li2025prompt} highlight emerging research directions including optimization of compression encoders through better architectural designs and more efficient information extraction, combining hard and soft methods to leverage complementary strengths, leveraging multimodality for cross-modal compression in image-text and audio-text scenarios, and advancing synthetic languages through more expressive learned tokens enabling universal compression languages with cross-model compatibility. These directions suggest that future compression methods will increasingly employ hybrid approaches that combine multiple compression strategies, adapt compression to task and content characteristics, and jointly optimize compression with downstream task performance.

The field's maturation is evident in the transition from isolated compression techniques to integrated systems that combine multiple methods. Production systems increasingly employ compression pipelines that chain complementary approaches, such as initial hard compression to remove obvious redundancy followed by soft compression to encode remaining content, or query-aware compression to identify relevant content followed by semantic compression of retained information. This compositional approach enables achieving compression ratios and quality levels unattainable by individual methods, suggesting that future advances may derive as much from effective composition of existing techniques as from novel compression algorithms.
